\section{Реконструкция изображений из мозговой активности}

Реконструкция изображений из мозговой активности представляет собой более сложную постановку, чем классификация когнитивного состояния или распознавание ограниченного набора стимулов. В классификации модель должна выбрать одну из заранее заданных меток, тогда как при реконструкции требуется восстановить структуру выходного объекта: форму, расположение элементов, контуры или семантическое содержание. Для EEG такая задача осложняется низким пространственным разрешением, низким отношением сигнал/шум и выраженной вариативностью между испытуемыми.

\subsection{От декодирования к реконструкции}

В исследованиях мозгового декодирования различают несколько уровней выходного представления. Наиболее простой уровень связан с распознаванием класса стимула или состояния. Более сложный уровень предполагает восстановление признаков изображения: геометрической структуры, положения элементов, низкоуровневых визуальных характеристик или латентного семантического представления. Наконец, генеративная реконструкция направлена на синтез визуально правдоподобного изображения с использованием GAN, VAE, diffusion models или трансформерных декодеров.

Современный обзор EEG-to-output исследований показывает, что направление активно развивается для изображений, видео и аудио, однако остаётся ограниченным качеством EEG-данных, отсутствием единых бенчмарков и трудностями cross-subject переноса~\footcite{10.48550/arxiv.2412.19999}. Поэтому в EEG-реконструкции важно отделять визуально впечатляющий генеративный результат от проверяемого восстановления информации, действительно содержащейся в сигнале.

\subsection{Восприятие и визуальное воображение}

Существенное различие существует между реконструкцией по данным прямого зрительного восприятия (visual perception) и по данным визуального воображения или воспоминания (visual imagery / memory recall). При восприятии стимул непосредственно представлен в сенсорном потоке и вызывает более устойчивые зрительные ответы. При воображении стимул должен быть извлечён из памяти и внутренне воспроизведён, поэтому сигнал может быть менее синхронизирован с событием и сильнее зависеть от стратегии испытуемого.

Нейрофизиологические исследования показывают, что восприятие и визуальное воображение могут иметь частично общие представления, в частности в альфа-диапазоне~\footcite{10.1016/j.cub.2020.04.074}. При этом их временная динамика различается, что требует осторожности при переносе методов, разработанных для perceptual decoding, на imagery decoding~\footcite{10.7554/elife.33904}. Для настоящей работы это принципиально: EEG регистрируется не во время предъявления изображения, а во время воспоминания ранее увиденного стимула.

\subsection{Семантическая и пиксельная реконструкция}

Многие современные подходы к реконструкции изображений используют семантические латентные пространства: модель сначала предсказывает признаки, связанные с категорией или смысловым содержанием стимула, а затем генеративная модель синтезирует изображение. Такой подход может давать визуально правдоподобные результаты, но не всегда гарантирует точное восстановление конкретной геометрии исходного стимула.

Для задачи реконструкции бинарного изображения $6 \times 6$ пикселей более уместна локальная пиксельная постановка. В ней каждый элемент изображения рассматривается как отдельная целевая переменная или часть многовыходного предсказания. Такой подход не стремится генерировать фотореалистичное изображение, но позволяет проверять, насколько EEG-признаки несут информацию о пространственной структуре простого визуального стимула.

Пиксельная постановка также лучше согласуется с ограничениями EEG. Низкое пространственное разрешение и объёмная проводимость затрудняют восстановление сложных деталей, однако простые бинарные стимулы малого размера уменьшают размер выходного пространства и делают задачу экспериментально контролируемой. Это позволяет оценивать качество реконструкции с помощью понятных метрик, таких как pixel accuracy, F1-score и intersection over union (IoU).

\subsection{Метрики оценки качества реконструкции}

Выбор метрик оценки зависит от того, какой тип выходного представления восстанавливает модель. В задачах классификации EEG обычно используются accuracy, balanced accuracy, precision, recall, F1-score и площадь под ROC-кривой (area under the ROC curve, ROC-AUC). Эти показатели удобны, когда модель предсказывает дискретную метку состояния или класса стимула, однако они не всегда отражают качество восстановления пространственной структуры изображения.

Для пиксельной реконструкции бинарных изображений естественно рассматривать задачу как набор связанных бинарных решений. В этом случае применяются pixel accuracy, precision, recall и F1-score по пикселям. Pixel accuracy показывает долю правильно восстановленных элементов, но может быть завышена при дисбалансе между фоном и активными пикселями. F1-score лучше отражает баланс между ложноположительными и ложноотрицательными пикселями, особенно если активных элементов изображения мало.

Отдельную группу составляют метрики перекрытия, такие как intersection over union (IoU) и Dice coefficient. Они сравнивают область активных пикселей в истинном и реконструированном изображениях и поэтому лучше соответствуют задаче восстановления формы. Для бинарных стимулов малого размера такие метрики особенно информативны, поскольку позволяют оценить не только число совпавших пикселей, но и пространственное совпадение восстановленной структуры.

В задачах непрерывной или полутоновой реконструкции изображений часто используются среднеквадратичная ошибка (mean squared error, MSE), средняя абсолютная ошибка (mean absolute error, MAE), peak signal-to-noise ratio (PSNR) и structural similarity index measure (SSIM). Эти метрики оценивают различие интенсивностей и структурное сходство изображений. Однако для бинарного изображения $6 \times 6$ их применимость ограничена: MSE и MAE могут быть полезны при вероятностном выходе модели, тогда как SSIM и PSNR более естественны для изображений с богатой локальной структурой и большим разрешением.

Для генеративных EEG-to-image подходов также применяются перцептивные и семантические метрики, например learned perceptual image patch similarity (LPIPS), Fréchet inception distance (FID) и сходство в признаковом пространстве CLIP. Они оценивают визуальную правдоподобность и семантическое соответствие, но хуже подходят для строгой проверки восстановления конкретного бинарного стимула. В настоящей постановке приоритет имеют метрики, непосредственно связанные с совпадением пикселей и формы, а не с субъективной визуальной реалистичностью результата.

Наконец, для EEG-реконструкции важна не только сама метрика, но и способ её статистической интерпретации. Результаты следует сравнивать со случайным baseline, majority baseline и простыми моделями, а также анализировать отдельно по испытуемым и по протоколам intra-subject и inter-subject. Это снижает риск переоценки качества модели в условиях малого числа эпох и выраженной межсубъектной вариативности.

\subsection{Ограничения EEG-реконструкции}

Основными ограничениями EEG-реконструкции являются слабая локализация источников, шумность, артефакты, межсубъектная вариативность и ограниченное число размеченных эпох. Эти факторы повышают риск того, что модель будет восстанавливать не визуальное содержание, а побочные корреляции: особенности отдельного испытуемого, временной структуры протокола или артефактов регистрации.

Поэтому для корректной постановки необходимо использовать строгие схемы разделения данных, сравнивать результаты с простыми baseline-моделями и отдельно анализировать intra-subject и inter-subject протоколы. Если модель показывает качество выше случайного уровня только внутри одного испытуемого, это следует интерпретировать как индивидуально-специфическое декодирование, а не как универсальное восстановление визуального образа.

\subsection*{Выводы по разделу}

Реконструкция изображений из EEG является перспективным, но методологически сложным направлением. В отличие от задач классификации, она требует восстановления структуры выходного объекта, а в случае visual imagery дополнительно сталкивается с нестабильностью внутреннего воспроизведения образа. На этом фоне локальная реконструкция бинарных изображений $6 \times 6$ представляет собой обоснованный компромисс: она сохраняет связь с задачей восстановления визуальной структуры, но избегает чрезмерных требований к пространственной точности EEG и объёму обучающих данных. Для такой постановки наиболее релевантны pixel-level и shape-overlap метрики, дополняемые сравнением с baseline-моделями и раздельным анализом intra-subject и inter-subject сценариев.
